\subsection{多线性映射关于由半范数定义的拓扑的等度连续性判别法}

命题 4. --- 设 \(E_{i}\)（\(1 \leqslant i \leqslant n\)）与 F 是 K 上的拓扑向量空间；我们假设，对每个 i，\(E_{i}\) 的拓扑由一个有向半范数集 \(\Gamma_{i}\) 定义，而 F 的拓扑由半范数集 \(\Gamma\) 定义。那么，\(\prod_{i=1}^{n} \mathrm{E}_{i}\) 到 F 中的多线性映射所成的一个集 H 是等度连续的，当且仅当：对每个半范数 \(q \in \Gamma\) 及每个指标 \(i\)，存在一个半范数 \(p_{i} \in \Gamma_{i}\) 及一个数 \(a > 0\)，使得对每个映射 \(u \in \mathrm{H}\) 及每点 \((x_{i}) \in \prod_{i=1}^{n} \mathrm{E}_{i}\)，

\[
q \left(u \left(x _ {1}, x _ {2}, \dots , x _ {n}\right)\right) \leqslant a. p _ {1} \left(x _ {1}\right) p _ {2} \left(x _ {2}\right) \dots p _ {n} \left(x _ {n}\right).\tag{5}
\]

条件是充分的，因为它蕴含 H 在 \((0, 0, \ldots, 0)\) 处等度连续，从而处处等度连续（I, 第 9 页, prop. 6）。

我们证明条件是必要的。由假设，对每个半范数 \(q \in \Gamma\) 及每个数 \(\beta > 0\)，都有 \(q(u(x_1, x_2, ..., x_n)) \leqslant \beta\)（对每个映射 \(u \in \mathbf{H}\)），只要 \(p_i(x_i) \leqslant \alpha_i\) 对每个指标 \(i\)（\(1 \leqslant i \leqslant n\)）成立，而数 \(\alpha_i > 0\) 与半范数 \(p_i \in \Gamma_i\) 是适当选取的。由于 \(\mathbf{K}\) 非离散，我们还可假设对每个 \(i\)，\(\alpha_i = |\lambda_i| < 1\)，其中 \(\lambda_i \in \mathbf{K}\)。现在设 \((x_1, x_2, ..., x_n)\) 是 \(\prod_{i=1}^{n} E_i\) 中任意一点，并对每个指标 \(i\)，设 \(m_i \in \mathbf{Z}\) 是满足 \(p_i(x_i) \leqslant |\lambda_i|^{m_i + 1}\) 的整数；这可写成 \(p_i(\lambda_i^{-m_i} x_i) \leqslant |\lambda_i| (1 \leqslant i \leqslant n)\)，因此由假设，我们有

\[
q \left(u \left(x _ {1}, x _ {2}, \dots , x _ {n}\right)\right) \leqslant \beta \left| \lambda_ {1} \right| ^ {m _ {1}} \left| \lambda_ {2} \right| ^ {m _ {2}} \dots \left| \lambda_ {n} \right| ^ {m _ {n}}.\tag{6}
\]

首先假设某个 \(p_{i}(x_{i})\) 为零，那么可取 \(m_{i} \in N\) 任意大，因此 \(q(u(x_{1}, x_{2}, ..., x_{n})) = 0\)。反之，若所有 \(p_{i}(x_{i})\) 都 \(\neq 0\)，则对每个 i 取整数 \(m_{i}\) 使得 \(|\lambda_{i}|^{m_{i}+2} < p_{i}(x_{i}) \leqslant |\lambda_{i}|^{m_{i}+1}\)；于是有 \(|\lambda_{i}|^{m_{i}} < |\lambda_{i}|^{-2}p_{i}(x_{i})\)，由此并根据 (6)，即得关系式 (5)，其中

\[
a = \beta (| \lambda_ {1} | \cdot | \lambda_ {2} | \dots | \lambda_ {n} |) ^ {- 2}. \qquad \text{Q.E.D.}
\]

推论. --- 集 H 是等度连续的，当且仅当：对每个半范数 \(q \in \Gamma\)，存在 \(\prod_{i=1}^{n} E_{i}\) 中 0 的一个邻域，使得诸函数 \(q \circ u\)（\(u \in H\)）在其中一致有界。

条件显然是必要的，而且 prop. 4 的证明表明它蕴含对所有 \(u \in H\) 成立的形如 (5) 的不等式，从而蕴含 H 的等度连续性。

我们把 prop. 4 关于线性映射的特例明确叙述出来。

命题 5. --- 设 E、F 是非离散赋值除环 K 上的两个拓扑向量空间；假设 E（相应地 F）的拓扑由半范数集 \(\Gamma\)（相应地 \(\Gamma'\)）定义。设 H 是 E 到 F 中的线性映射组成的一个集。下列条件等价：

\begin{enumerate}
\def\labelenumi{\alph{enumi})}
\item
  H 等度连续。
\item
  对每个半范数 \(q \in \Gamma'\)，存在半范数的一个有限族 \((p_i)_{1 \leq i \leq n}\)（属于 \(\Gamma\)）及一个数 \(a > 0\)，使得对所有 \(x \in E\) 及所有 \(u \in H\)，
\end{enumerate}

\[
q (u (x)) \leqslant a. \sup _ {1 \leqslant i \leqslant n} p _ {i} (x).\tag{7}
\]

\begin{enumerate}
\def\labelenumi{\alph{enumi})}
\setcounter{enumi}{2}
\tightlist
\item
  对每个半范数 \(q \in \Gamma'\)，映射 \(\sup_{u \in H} (q \circ u)\) 是 \(E\) 上的一个连续半范数。
\end{enumerate}

推论 1. --- 设 \(\mathcal{T}, \mathcal{T}'\) 是向量空间 \(\mathbf{E}\)（在 \(\mathbf{K}\) 上）上分别由两个半范数集 \(\Gamma\) 与 \(\Gamma'\) 定义的拓扑。\(\mathcal{T}\) 比 \(\mathcal{T}'\) 更细，当且仅当：对每个半范数 \(q \in \Gamma'\)，存在半范数的一个有限族 \((p_i)_{1 \leqslant i \leqslant n}\)（属于 \(\Gamma\)）及一个数 \(a > 0\)，使得对所有 \(x \in \mathbf{E}\)，有 \(q(x) \leqslant a . \sup_{1 \leqslant i \leqslant n} p_i(x)\)。

事实上，这表明赋有拓扑 T 的 E 到赋有拓扑 \(T'\) 的 E 上的恒等映射连续。

推论 2. --- 假设 \(\mathcal{T}\) 是拓扑向量空间 \(E\)（在 \(K\) 上的）的拓扑，由一个有向半范数集 \(\Gamma\) 定义；对每个半范数 \(p \in \Gamma\)，设 \(E_p\) 是由 \(E\) 赋有由 \(p\) 定义的拓扑而得的空间。集 \(E'\)（由 \(E\) 上关于 \(\mathcal{T}\) 连续的线性型组成）是诸集 \(E_p'\) 的并，其中 \(E_p'\) 是 \(E_p\) 中连续线性型组成的集（\(p \in \Gamma\)）。

